\subsection{Acoustic and electromagnetic guidance}

Acoustic-based navigation is widely used as the long-range guidance channel for underwater docking. A common configuration has an acoustic transceiver on the dock that emits a signal, which is received by a transponder on the AUV. The AUV then replies with its own signal, allowing the dock to compute the distance based on the time delay. To obtain bearing information, several inverted ultra-short baseline (iUSBL) receivers are spaced precisely on the AUV hull. By measuring the slight time difference of arrival of the acoustic signals at each receiver, the system can triangulate the relative azimuth and elevation angles between DS and AUV. WHOI's REMUS AUV docking system is a prime example of a USBL-navigation aided system \cite{4099107}. Its acoustic homing can guide the AUV from long range, at the cost of moderate update rates, latency due to the speed of sound underwater, risk of false signals from acoustic refraction and ambient noise, and limited angular resolution.

Electromagnetic (EM) guidance offers an alternative to acoustic signals, using low-frequency magnetic fields emitted by a coil on the dock, and results in lower latency. Feezor et al. (2001) developed an EM homing system where a low-frequency magnetic field is generated by a coil on the dock, and the AUV uses magnetometers to detect the field strength and direction, without line-of-sight constraints \cite{feezor2001autonomous}. Their experiments demonstrated that the AUV could home in on the dock from a limited range up to $\sim$30~m. Further, the system requires dedicated coils, drivers and calibration of the local field environment, which may limit its practicality for general-purpose docking stations.

Most underwater docking stations incorporate mechanical alignment aids to relax the sensing requirement and accommodate for the slight tolerances during the final docking alignment phase. There exist two main configurations: unidirectional and omnidirectional guides. Unidirectional stations are the most common type of configuration, typically designed for a torpedo-shaped vehicle approaching on a known heading. Omnidirectional stations allow approach from any heading, typically as a vertical pole or tensioned cable, allowing the vehicle to be latched onto the dock despite its orientation, notably in the Odyssey IIb AUV's dock system \cite{972084}, simplifying the navigation and alignment process at the cost of a dedicated latch mechanism on the vehicle nose. In either configuration, the guide's capture radius and tolerated closing speed define the interface between the guidance problem and the mechanical one, in which the capture envelope assumed in this work (Section~\ref{sec:method-experiment}) plays that role.
